% =================================================================
\section{Conclusion}
\label{sec:conclusion}
% =================================================================

The main lesson of this work is simple to state but easy to miss in
practice: if your training data is dominated by one condition, your
domain-generalisation method is not working---it is just learning
that one condition.
Fixing the imbalance first is not optional.
In the CD-MSC dataset, where one recording environment contributes
99.4\% of training clips, resampling to balance domains gains +15.6~pp
in cross-condition accuracy before any DG method is applied.
On top of that foundation, domain-adversarial training and
cross-condition contrastive learning push the system to $\BAu{=}0.378$
(+21.3~pp total) with a domain gap of $\DSG{=}0.202$.
FBS-Mix---which mixes only the recording-noise frequency band and
protects the wingbeat signal---provides a lightweight alternative that
reduces the seen/unseen gap to 0.229 with no model changes.

We also find that gradient reversal does not work for transformer
backbones in this setting: AST's domain classifier maintains 99.7\%
accuracy throughout training, meaning the adversarial signal never
suppresses condition information.
This failure comes from two compounding causes: the much larger
backbone gradient overwhelms the adversarial signal, and
per-sample layer normalisation removes the batch-level condition
statistics that gradient reversal relies on.

\paragraph{Limitations.}
All ablations use a single seed; full multi-seed variance estimates
for the proposed system are pending.
Domain labels are opaque, so we cannot say what physical differences
(device type, location, recording protocol) drive the cross-domain
gap.

\paragraph{Future work.}
The most promising direction for AST is applying domain-adversarial
training at the patch-embedding level rather than the output-token
level, where condition information may be more concentrated, or
using pretrained AST weights where some domain invariance is already
encoded.
For the MTRCNN system, combining FBS-Mix with the DiCL projection
head~(currently pending as ``best + augmentations'' in
Table~\ref{tab:ablation}) should reduce the seen/unseen gap below
0.200 while maintaining high $\BAu$.
